\chapter{Methodik}
\label{ch:methodik}

\section{Versuchssude}
Gebraut wurden zwölf identische Sude eines Pale Ale (Stammwürze \SI{12.5}{\degree P}),
aufgeteilt auf vier Temperaturstufen mit je drei Wiederholungen.

\begin{table}[htbp]
  \centering
  \caption{Rezeptur und Gärparameter (pro \SI{20}{\liter} Sud)}
  \label{tab:rezept}
  \begin{tabular}{lrl}
    \toprule
    Zutat / Parameter & Menge   & Einheit \\
    \midrule
    Pale-Ale-Malz     & 4{,}2   & kg \\
    Cascade-Hopfen    & 60      & g \\
    Obergärige Hefe   & 11{,}5  & g \\
    Gärtemperatur     & 16 -- 24 & \si{\celsius} \\
    Gärdauer          & 7       & Tage \\
    \bottomrule
  \end{tabular}
\end{table}

\section{Analytik und Verkostung}
Die Ester wurden per Gaschromatografie bestimmt. Zusätzlich bewerteten 24 Personen
die Biere blind nach Fruchtigkeit, Sauberkeit und Gesamteindruck
(Skala 1--9), siehe \cref{tab:rezept} für die Rahmenbedingungen.
